\documentclass[11pt,a4paper]{article}
\usepackage[utf8]{inputenc}
\usepackage[T1]{fontenc}
\usepackage{tgheros}
\renewcommand{\familydefault}{\sfdefault}
\usepackage[varqu,scaled=0.92]{zi4}
\usepackage{microtype}
\DisableLigatures[-,<,>]{encoding=T1,family=tt*}
\usepackage{needspace}
\usepackage[a4paper,top=2.5cm,bottom=2.5cm,left=2.5cm,right=2.2cm]{geometry}
\usepackage{graphicx}
\graphicspath{{../../template/}{figures/}}
\usepackage{xcolor}
\usepackage{listings}
\usepackage{upquote}
\usepackage{titlesec}
\usepackage{tocloft}
\usepackage{enumitem}
\usepackage{booktabs}
\usepackage{array}
\usepackage{float}
\usepackage[hidelinks]{hyperref}

\setlength{\parindent}{0pt}
\setlength{\parskip}{0.55em}
\setlist{itemsep=0.3em,topsep=0.3em}
\pagestyle{plain}

\titleformat{\section}{\Large\bfseries}{\thesection}{0.9em}{}
\titleformat{\subsection}{\large\bfseries}{\thesubsection}{0.9em}{}
\titlespacing*{\section}{0pt}{1.6em}{0.7em}
\titlespacing*{\subsection}{0pt}{1.2em}{0.5em}

\renewcommand{\contentsname}{Cuprins}
\renewcommand{\cfttoctitlefont}{\LARGE\bfseries}
\renewcommand{\cftsecfont}{\bfseries}
\renewcommand{\cftsecpagefont}{\bfseries}
\renewcommand{\cftsubsecfont}{\small}
\renewcommand{\cftsubsecpagefont}{\small}
\setlength{\cftbeforesecskip}{0.6em}

\definecolor{codebg}{HTML}{F6F6F6}
\definecolor{coderule}{HTML}{8C8C8C}
\definecolor{codenum}{HTML}{8C8C8C}
\definecolor{codecomment}{HTML}{6A6A6A}

\lstdefinestyle{term}{
  basicstyle=\ttfamily\small,
  backgroundcolor=\color{codebg},
  frame=single,
  rulecolor=\color{coderule},
  framesep=4pt,
  numbers=left,
  numberstyle=\tiny\color{codenum},
  numbersep=9pt,
  xleftmargin=1.4em,
  breaklines=true,
  postbreak=\mbox{\textcolor{codenum}{$\hookrightarrow$}\space},
  columns=fullflexible,
  keepspaces=true,
  showstringspaces=false,
  upquote=true,
  aboveskip=0.9em,
  belowskip=0.9em
}
\lstdefinestyle{ccode}{
  style=term,
  language=C,
  keywordstyle=\bfseries,
  commentstyle=\color{codecomment},
  morekeywords={uint,uint64}
}
\lstset{style=term,captionpos=b}

\usepackage{caption}
\captionsetup[lstlisting]{labelsep=endash,justification=centering,singlelinecheck=false,font=small,skip=7pt}
\captionsetup[table]{labelsep=endash,font=small,skip=6pt}
\renewcommand{\lstlistingname}{Figure}
\AtBeginDocument{\counterwithin{lstlisting}{section}\counterwithin{table}{section}}

\newcommand{\plainsection}[1]{\section*{#1}\addcontentsline{toc}{section}{#1}}
\newcommand{\code}[1]{\texttt{#1}}

\newcommand{\utmtitlepage}{%
\begin{titlepage}
  \centering
  \includegraphics[width=4.3cm]{utm-logo.jpeg}\par
  \vspace{1.3cm}
  {\large\bfseries UNIVERSITATEA TEHNICĂ A MOLDOVEI\par}
  \vspace{0.25cm}
  Facultatea Calculatoare, Informatică și Microelectronică\par
  Departamentul Ingineria Software și Automatică\par
  \vspace{2.6cm}
  {\Large\bfseries LUCRARE DE LABORATOR NR.~\labnumber\par}
  \vspace{0.2cm}
  {\LARGE\bfseries \labtitle\par}
  \vspace{0.45cm}
  Disciplina: Criptografie și Securitate\par
  \ifx\variant\empty\else\vspace{0.2cm}Varianta \variant\par\fi
  \vspace{2.8cm}
  \begin{minipage}[t]{0.46\textwidth}
    \raggedright
    \textbf{A efectuat:}\\
    student gr. \studentgroup\\
    \studentname
  \end{minipage}\hfill
  \begin{minipage}[t]{0.46\textwidth}
    \raggedleft
    \textbf{A verificat:}\\
    conf. univ., dr. Aureliu \textsc{Zgureanu}
  \end{minipage}
  \vfill
  Chișinău — 2026\par
\end{titlepage}}

\usepackage{textcomp}
\usepackage{framed}
\newenvironment{cipherbox}{%
  \def\FrameCommand{\fboxsep=6pt\fboxrule=0.4pt\fcolorbox{coderule}{codebg}}%
  \MakeFramed{\advance\hsize-\width\FrameRestore}%
  \ttfamily\footnotesize\raggedright\frenchspacing\setlength{\parskip}{0pt}}%
  {\par\endMakeFramed}

\newcommand{\labnumber}{3}
\newcommand{\labtitle}{Polyalphabetic Ciphers: Vigenère}
\newcommand{\studentname}{Daniel Nicorici}
\newcommand{\studentgroup}{FAF-243}
\newcommand{\variant}{16}
\lstset{inputencoding=utf8,extendedchars=true,literate=
  {Ă}{{\u{A}}}1 {Â}{{\^{A}}}1 {Î}{{\^{I}}}1 {Ș}{{\textcommabelow{S}}}1 {Ț}{{\textcommabelow{T}}}1
  {ă}{{\u{a}}}1 {â}{{\^{a}}}1 {î}{{\^{i}}}1 {ș}{{\textcommabelow{s}}}1 {ț}{{\textcommabelow{t}}}1
  {Ş}{{\c{S}}}1 {Ţ}{{\c{T}}}1 {ş}{{\c{s}}}1 {ţ}{{\c{t}}}1 {è}{{\`{e}}}1 {é}{{\'{e}}}1}
\lstdefinestyle{py}{style=term,language=Python,keywordstyle=\bfseries,commentstyle=\color{codecomment}}
\begin{document}
\utmtitlepage
\tableofcontents
\newpage

\plainsection{Purpose of the work}

A monoalphabetic cipher keeps the frequency of every letter, which is why Lab 2 could break it by
counting. A polyalphabetic cipher changes the substitution from one position to the next, so the
same plaintext letter is encrypted into different cipher letters and the frequencies flatten out.
The purpose of this lab is to implement the best known polyalphabetic cipher, the Vigenère cipher,
for messages in Romanian.

My number in the group register is 16, which is even, so I solved task 3.2: a Vigenère cipher over
the 31-letter Romanian alphabet, with letters coded 0 to 30. The program accepts only Romanian
letters and tells the user which characters are allowed, refuses keys shorter than 7 letters,
removes spaces and converts everything to upper case before encrypting, and lets the user choose
between encryption and decryption. It is written in Python 3 and needs no external libraries.

\section{The Vigenère cipher over 31 letters}

The Romanian alphabet adds five letters with diacritics to the 26 Latin ones. Each letter gets a
number from 0 to 30 in alphabetical order (Table~\ref{tab:alpha}).

\begin{table}[H]
\centering
\caption{The Romanian alphabet and the code of every letter}\label{tab:alpha}
\setlength{\tabcolsep}{3.6pt}
\small
\begin{tabular}{l*{16}{c}}
\toprule
Letter & A & Ă & Â & B & C & D & E & F & G & H & I & Î & J & K & L & M \\
Code & 0 & 1 & 2 & 3 & 4 & 5 & 6 & 7 & 8 & 9 & 10 & 11 & 12 & 13 & 14 & 15 \\
\midrule
Letter & N & O & P & Q & R & S & Ș & T & Ț & U & V & W & X & Y & Z & \\
Code & 16 & 17 & 18 & 19 & 20 & 21 & 22 & 23 & 24 & 25 & 26 & 27 & 28 & 29 & 30 & \\
\bottomrule
\end{tabular}
\end{table}

The key is a word $k_1 k_2 \dots k_t$ of at least 7 letters. It is written under the message and
repeated as many times as needed, so the letter at position $i$ is shifted by the key letter
$k_{(i \bmod t)}$. The mathematical model from the handout becomes, for 31 letters:
\[
  c_i = (m_i + k_i) \bmod 31, \qquad m_i = (c_i - k_i) \bmod 31 .
\]

\textbf{Observations:}
\begin{enumerate}
  \item \textbf{Every key letter is a Caesar cipher.} With a key of length $t$, positions $1, t+1,
  2t+1, \dots$ are all shifted by $k_1$, positions $2, t+2, \dots$ by $k_2$, and so on. Vigenère
  is $t$ interleaved Caesar ciphers, which is why it uses the same operation as Caesar.
  \item \textbf{The key space grows with the key.} There are $31^t$ keys of length $t$. The
  shortest allowed key, 7 letters, already gives $31^7 \approx 2.75 \cdot 10^{10}$ keys, and the
  12-letter key used below gives $31^{12} \approx 7.9 \cdot 10^{17}$.
  \item \textbf{Modulo 31 instead of 26.} The formulas stay the same; only the modulus and the
  table of codes change. A shift of 1 now turns \code{Z} (30) into \code{A} (0), and \code{B} has
  code 3, not 1, because \code{Ă} and \code{Â} come before it.
\end{enumerate}

\section{Implementation}

The whole cipher is one function, \code{shift}, which adds or subtracts the key codes depending on
its \code{sign} argument; \code{encrypt} calls it with $+1$ and \code{decrypt} with $-1$, exactly
as in the two formulas above. The other functions prepare and check the input.

\begin{itemize}
  \item \code{clean} removes every space (and tab or newline), converts the text to upper case and
  checks that only letters of the alphabet remain. Any other character raises an error that lists
  the wrong characters and names the allowed range. The old cedilla forms \code{Ş} and \code{Ţ},
  which many keyboards still produce, are first replaced by the correct comma-below letters
  \code{Ș} and \code{Ț}, so a message typed either way is accepted.
  \item \code{check\_key} applies the same rules to the key and also refuses keys shorter than 7
  letters.
  \item \code{ask} repeats a question until the answer passes its check, so a wrong input never
  stops the program.
  \item \code{main} shows the menu: 1 encrypts, 2 decrypts, 0 exits.
\end{itemize}

The alphabet is a parameter of \code{encrypt} and \code{decrypt}, with the Romanian alphabet as the
default. This let me check the program against the 26-letter example of the handout
(section~\ref{sec:tests}) without a second implementation.

\lstinputlisting[style=py,caption={vigenere.py},label={fig:code}]{../vigenere.py}

\section{Running the program}

The session below first tries a key that is too short and a message with a digit and punctuation,
both of which are refused, then encrypts a Romanian sentence with the key \code{criptografie} and
decrypts the result again. The message is typed with cedilla letters (\code{ş}, \code{ţ}) and in
lower case on purpose.

\begin{lstlisting}[caption={A session with the program},label={fig:run}]
$ python3 vigenere.py
Cifrul Vigenère pentru alfabetul român (31 de litere, codificate 0..30)

1 - criptare, 2 - decriptare, 0 - ieșire: 1
Cheia: cheie
Cheia are 5 litere. Lungimea cheii trebuie să fie de cel puțin 7 litere.
Cheia: criptografie
Mesajul: Ora 7, la poartă!
Caractere nepermise: ! , 7. Sunt permise doar literele alfabetului român: A-Z, a-z și Ă Â Î Ș Ț, ă â î ș ț.
Mesajul: Întâlnirea secretă are loc mâine la ora şapte lângă ţărmul lacului
Mesajul pregătit: ÎNTÂLNIREASECRETĂARELOCMÂINELAORAȘAPTELÂNGĂȚĂRMULLACULUI
Criptograma:      MDÂREÂPHEFAJGHNIȚOXVLȚLSEZVȚEOUHAYIȚWVȚRGUHKĂWUAPBIȘOAÂZ

1 - criptare, 2 - decriptare, 0 - ieșire: 2
Cheia: criptografie
Criptograma: MDÂREÂPHEFAJGHNIȚOXVLȚLSEZVȚEOUHAYIȚWVȚRGUHKĂWUAPBIȘOAÂZ
Mesajul decriptat: ÎNTÂLNIREASECRETĂARELOCMÂINELAORAȘAPTELÂNGĂȚĂRMULLACULUI

1 - criptare, 2 - decriptare, 0 - ieșire: 3
Alegeți 1, 2 sau 0.

1 - criptare, 2 - decriptare, 0 - ieșire: 0
\end{lstlisting}

To show what the program computes, Table~\ref{tab:hand} repeats the encryption of the first 12
letters by hand, the same way as the table in the handout. Twelve letters are exactly one period
of the key.

\begin{table}[H]
\centering
\caption{Encryption of the first 12 letters with the key CRIPTOGRAFIE}\label{tab:hand}
\setlength{\tabcolsep}{4.2pt}
\small
\begin{tabular}{l*{12}{c}}
\toprule
Plaintext $m$ & Î & N & T & Â & L & N & I & R & E & A & S & E \\
Code $m_i$ & 11 & 16 & 23 & 2 & 14 & 16 & 10 & 20 & 6 & 0 & 21 & 6 \\
Key $k$ & C & R & I & P & T & O & G & R & A & F & I & E \\
Code $k_i$ & 4 & 20 & 10 & 18 & 23 & 17 & 8 & 20 & 0 & 7 & 10 & 6 \\
$(m_i + k_i) \bmod 31$ & 15 & 5 & 2 & 20 & 6 & 2 & 18 & 9 & 6 & 7 & 0 & 12 \\
Ciphertext $c$ & M & D & Â & R & E & Â & P & H & E & F & A & J \\
\bottomrule
\end{tabular}
\end{table}

\textbf{Observations:}
\begin{enumerate}
  \item \textbf{The same letter is encrypted differently.} \code{N} at positions 2 and 6 becomes
  \code{D} and then \code{Â}, and \code{E} at positions 9 and 12 becomes \code{E} and then
  \code{J}. In the whole message the 6 letters \code{E} turn into 5 different cipher letters
  (\code{E}, \code{J}, \code{N}, \code{V}, \code{Ț}) and the 6 letters \code{A} into 4
  (\code{A}, \code{F}, \code{I}, \code{O}). This is the point of a polyalphabetic cipher: the most
  frequent plaintext letter no longer produces one most frequent cipher letter, so the attack of
  Lab 2 does not work.
  \item \textbf{Different letters can give the same cipher letter.} \code{T} and \code{N} at
  positions 3 and 6 both become \code{Â}, because $23 + 10 = 33 \equiv 2$ and
  $16 + 17 = 33 \equiv 2 \pmod{31}$. High and low frequencies get mixed, which is the flattening
  described at the start of the handout.
  \item \textbf{A key letter A changes nothing.} Its code is 0, so at position 9 the letter
  \code{E} stays \code{E}. A key full of A's would leave the message readable, which is one more
  reason for a long key without repeated patterns.
  \item \textbf{Spaces are gone for good.} The decrypted text has no spaces, so, as the task
  says, splitting it back into words is done by the reader.
\end{enumerate}

\section{Testing}\label{sec:tests}

The test file checks the program against the handout and against itself: the 26-letter example
\textit{Per aspera ad astra} with key \code{SUPER} must give \code{HYGEJHYGERVUHXIS}, decryption
must undo encryption for a Romanian message, the wrap-around from \code{Z} to \code{A} must work,
characters outside the alphabet must be refused, and a 5-letter key must be refused.

\lstinputlisting[style=py,caption={test\_vigenere.py},label={fig:tests}]{../test_vigenere.py}

\begin{lstlisting}[caption={Running the tests},label={fig:testrun}]
$ python3 test_vigenere.py
all tests passed
\end{lstlisting}

\plainsection{Conclusions}

The Vigenère cipher was implemented for the 31-letter Romanian alphabet in about seventy lines of
Python. Encryption and decryption are the same function with the sign of the key reversed, which
follows directly from the formulas $c_i = (m_i + k_i) \bmod 31$ and $m_i = (c_i - k_i) \bmod 31$.
The program reproduces the example from the handout when given the 26-letter alphabet, and every
Romanian message tested decrypts back to itself.

Most of the work was in preparing the input rather than in the cipher. Spaces have to be removed,
lower case converted, the cedilla letters Ş and Ţ mapped to the correct Ș and Ț, and every other
character refused with a message that tells the user what is allowed.

The example shows why the cipher resists the frequency analysis used in Lab 2: one plaintext letter
is turned into several cipher letters and several plaintext letters share one cipher letter, so the
counts no longer point at e or a. It is still not secure today: the key repeats, so once the key
length is found (Kasiski examination, index of coincidence) the cryptogram splits into a few Caesar
ciphers, each broken by counting letters again.
\end{document}
